\noindent Desde que Kurt Gödel demostró sus famosos teoremas de incompletitud en 1931, los matemáticos se han encontrado con una gran cantidad de enunciados que, aunque pueden ser escritos en una cierta teoría, no pueden ser demostrados en la misma.
En 1944, Reuben Goodstein demostró que las sucesiones que a día de hoy llevan su nombre, un tipo de sucesiones que crecen muy rápidamente, no sólo no divergen, sino que convergen todas y a 0. Para lograr demostrar esto se valió de la, aún joven, teoría de conjuntos y  dio una prueba bastante concreta con la que quedó bien claro que las sucesiones de Goodstein convergen a 0. Sin embargo, observemos un momento el enunciado del teorema:
\[(\forall m)(\exists k)(g_{k}(m)=0)\]
donde $g_k(m)$ representa al $k-$ésimo elemento de la sucesión de Goodstein asociada al número $m$.
Si bien no se alcanza a apreciar que la definición de sucesión de Goodstein sea posible en la aritmética de Peano, ciertamente lo es (y se aborda más adelante en este trabajo), por lo que podemos notar que el enunciado del teorema también se puede escribir perfectamente sin salirnos de la aritmética de Peano, la cuál estudiaremos y definiremos ampliamente en el capítulo {\ref{axiomas_peano}} de este mismo trabajo.
De aquí, surge la pregunta natural: ¿es posible demostrar el Teorema de Goodstein utilizando únicamente como base los axiomas que Giuseppe Peano publicó en 1889 (o un conjunto equivalente de axiomas)?
La respuesta la dieron en 1982 Laurence Kirby y Jeff Paris en 1982 en el trabajo \textit{Accessible independence results for Peano arithmetic}{\cite{Goodstein}} y dicha respuesta era un rotundo \textit{no}.

En este trabajo abordaremos como eje central el \textit{teorema de Goodstein}, un teorema que habla de unas ciertas sucesiones de números naturales junto con su extraño y completamente antiintuitivo comportamiento. Además, estudiamos de forma paralela unas extrañas estructuras llamadas \textit{hidras}, las cuales tienen un comportamiento parecido a las sucesiones de Goodstein.
Para esto nos hemos valido de la teoría de conjuntos, en especial de la teoría de ordinales; así como de herramientas de la teoría de modelos. Todo esto con el objetivo de dar una prueba sólida, y lo mejor explicada posible, sobre la imposibilidad de demostrar el teorema de Goodstein bajo la aritmética de Peano.
Asimismo, desarrollaremos las herramientas necesarias para encontrar las conexiones entre las sucesiones de Goodstein y las hidras, y también demostraremos, en el capitulo \ref{indemostrabilidad}, que el hecho de que las hidras puedan ser derrotadas bajo cualquier estrategia es algo indemostrable bajo los mismos axiomas de Peano. 
Ambas pruebas de indemostrabilidad las daremos por medio de teoremas distintos, aunque vale la pena aclarar que, dadas las conexiones, ambos resultados se podrían resolver de la misma forma.
La motivación de este trabajo es dar una demostración lo más completa posible, pero que a su vez sea entendible, pues en el ya mencionado de Kirby y Paris se da por hecho que el lector tiene un barage amplio sobre conocimientos dentro de la teoría de conjuntos, teoría de la computabilidad, teoría de modelos, etcétera.

En el primer capítulo, nos encargamos de dar todo resultado preliminar que se podría abordar en un curso de licenciatura en matemáticas o equivalente.
En el segundo capítulo presentamos nuestros objetos de estudio: las sucesiones de Goodstein y las hidras
En el tercer capítulo desarrollamos nuestra teoría de la aritmética, es decir, aquellas consecuencias de los axiomas de Peano, y también veremos la herramienta más importante de todo este trabajo: los conjuntos $\alpha-$grandes.
En el cuarto capítulo veremos las últimas herramientas necesarias para probar la indemostrabilidad de los teoremas mencionados.
Finalmente en el apéndice A mencionaremos una equivalencia útil en muchas áreas de las matemáticas, y daremos una prueba de que las sucesiones de Goodstein son funciones computables.

Así pues, los prerrequisitos que se piden para leer este trabajo es un conocimiento básico en teoría de conjuntos. En el primer capítulo introduciremos a los ordinales, una de las herramientas matemáticas más poderosas que introdujo el mismísimo Georg Cantor en 1897, explicaremos su definición, sus propiedades básicas, su relación con los conjuntos bien ordenados, cómo dan paso a uno de los tipos de demostración más fuerte que existen (la inducción transfinita) y cómo operar con ellos mediante reglas aritméticas.
También abordaremos la herramienta de los ultrafiltros, un tipo de familia de conjuntos que nos va a permitir más adelante hablar sobre modelos no estándar dentro de nuestras teorías matemáticas. Los metemos acá porque es un concepto no muy complicado que es normalmente abordado en un curso usual de Topología en licenciatura.
Asimismo, introduciremos al lector en temas de lógica matemática, especialmente en lógica de primer orden, mencionando el cálculo deductivo y las bases detrás de lo que son los teoremas, los modelos y las deducciones en matemáticas.
Por otro lado, daremos una introducción a la teoría de la computabilidad, teoría introducida en la decada de 1930 y cuyo principal representante es Alan Turing, esto con el fin de abordar temas como las máquinas de Turing, las funciones computables y los conjuntos computables, todos temas que serán de capital importancia en el cuarto capítulo de este trabajo.
Finalmente, y por completud, mencionaremos un espacio para los teoremas de incompletitud de Kurt Gödel, pues todo lo expuesto anteriormente permite ya abordar temas así de importantes y que sería un desperdicio no mencionar.

El segundo capítulo funciona como introducción a los dos conceptos principales de nuestro trabajo: las sucesiones de Goodstein y las hidras. Este capítulo es corto, pues su función es presentar los conceptos y dar las demostraciones (en teoría de conjuntos) de que ambos conceptos pueden ser derrotados.\footnote{El concepto de derrotar una hidra hace referencia a cortar todas las cabezas de esta misma y es un tema que explicaremos ampliamente en el capítulo 2.2.} También sirve para explicar cosas que más adelante daremos completamente por entendido, como la formalización en aritmética de Peano de las sucesiones y la asociación con ordinales de las hidras.

En el tercer capítulo, primero presentaremos el sistema axiomático en el que estamos trabajando, los multinombrados axiomas de Peano. Esto con el fin de que después, en este mismo capítulo, podamos hablar sobre los ultraproductos, un concepto que se basa en los ultrafiltros del primer cápitulo, y que son la base para dar con un modelo no estándar de la aritmética de Peano, es decir, un conjunto de números que se va a portar como los números naturales (pues sigue los axiomas de Peano), pero que al no ser los números naturales va a arrojarnos resultados extraños, que permitirán demostrar que los axiomas de Peano no implican el cumplimiento del teorema de Goodstein.
Por otro lado, en este mismo capítulo introduciremos a los conjuntos $\alpha-$grandes, una medida para saber qué tan grande es un intervalo y, de esa forma, estudiar como se comportarán las sucesiones de Goodstein dentro de dicho intervalo.

En el capítulo cuatro abordaremos en profundidad herramientas mencionadas o explicadas en el trabajo de Paris y Kirby, que ya mencionamos anteriormente. Proveeremos pruebas hacia los teoremas de Ketonen-Solovay, realizaremos cálculos y abordaremos conceptos de la teoría de computabilidad que en el artículo original se daban simplemente por citadas, todo para tratar de dejar todo claro, coherente y completo.
Finalmente, en este capítulo, daremos la demostración de que tanto la convergencia de las sucesiones de Goodstein, como la derrota de las hidras, son teoremas que no se pueden demostrar únicamente usando las herramientas de la aritmética de Peano.

En el apéndice A, mencionamos un hecho super conocido sobre una de las equivalencias más conocidas dentro de las matemáticas: el axioma de elección es equivalente al lema de Zorn. 
Y finalmente, en el apéndice B, daremos de forma explícita un código en Python para calcular sucesiones de Goodstein, lo cuál sería un demostración sobre que las sucesiones de Goodstein son objetos computables.
De esta forma, quizá valga la pena resaltar, también damos un ejemplo de que computabilidad no implica demostrabilidad.
